\section{Introduction}
\label{sec:intro}

HPC systems are becoming interactive scientific environments rather than purely batch-oriented ones. Users increasingly delegate routine but consequential work --- reading logs, repairing job scripts, installing software, summarizing large output directories, and steering multi-stage simulations --- to LLM agents that can read files, invoke command-line tools, call scheduler commands, and communicate with local or remote model backends. This is not speculative: agent frameworks have already been connected to real HPC resource managers and tested on production-scale scientific workflows \cite{Ma2025HPCAgent,Rosendo2025ORNL,Dawson2026LARA}, multi-agent systems have been used to automate HPC software compilation and deployment \cite{Mondesire2025PEARC}, and HPC centers are already publishing user guidance about agentic coding tools appearing on their clusters \cite{AaltoSciComp2026}. The trajectory is toward more autonomy, not less.

This creates a security problem that does not fit cleanly into either of the two bodies of work that study its component parts.

LLM-agent security research has established that agents which consume untrusted external content can be manipulated through indirect prompt injection \cite{Greshake2023PromptInjection}, that tool-using agents fail in high-stakes ways even when reasonably careful \cite{Ruan2024ToolEmu}, that these failures can be measured with realistic tasks and adaptive attacks \cite{Debenedetti2024AgentDojo}, and that the attack surface extends to tool selection \cite{Shi2025ToolHijacker} and inter-agent/inter-protocol communication \cite{Ferrag2025ProtocolExploits,Wang2026PILandscape}. But the settings in which this is studied --- email clients, e-banking, travel booking, Slack, generic tool sandboxes --- do not have schedulers, shared parallel filesystems, project-scoped multi-tenant storage, or the notion that a valid user account may span multiple projects while a given task should not.

HPC security research has, separately, built mature answers to ``who can access what'': zone-based reference architectures for access, management, compute, and storage \cite{NIST2024SP800223}; enforced per-user separation of processes, filesystem access, network traffic, and accelerators \cite{Prout2024UserSeparation}; and federated identity and zero-trust designs that unify authentication and access control across AI and HPC infrastructure \cite{Alam2024ZeroTrust}. All of this assumes that the entity making a decision behind an authenticated session is a human, or software the human directly and deterministically controls. None of it asks whether a \emph{sequence of authorized actions}, taken by software that reads adversarial content as part of its normal operation, correctly reflects what the user or the site intended.

The hijacked authorized agent threat model exists in the space between these two literatures: \textbf{the identity is real, the credentials are valid, the authorization is technically correct --- and the behavior is not what anyone intended.}

\subsection{Thesis}
\label{sec:thesis}

We argue three things in this paper:

\begin{enumerate}[leftmargin=*]
  \item \textbf{The threat is structurally different from generic prompt injection}, because HPC introduces authority and persistence dimensions --- scheduler-mediated resource consumption, multi-project account scope, shared and durable parallel-filesystem state, and scientific-integrity/provenance requirements --- that general agent-security settings do not have and that determine what ``unsafe'' even means in this domain.
  \item \textbf{Neither literature currently measures it.} LLM-agent security benchmarks do not model scheduler authority or shared HPC storage; HPC security controls do not model semantic redirection of an authenticated agent.
  \item \textbf{The gap is closing on its own schedule, not ours.} HPC agents are being deployed for real workflows now \cite{Ma2025HPCAgent,Rosendo2025ORNL,Mondesire2025PEARC,AaltoSciComp2026}, and site operators are already worried enough to write informal guidance about it \cite{AaltoSciComp2026}, which means the research community has a limited window to characterize the threat before ad hoc, under-evaluated mitigations become the de facto standard.
\end{enumerate}

\subsection{Contributions}
\label{sec:contributions}

\begin{enumerate}[leftmargin=*]
  \item A precise definition of the \textbf{hijacked authorized agent} threat model for HPC, distinguishing it from conventional privilege escalation and from generic indirect prompt injection (Section~\ref{sec:threat-model}).
  \item A \textbf{taxonomy of HPC-specific attack surfaces}: shared parallel-filesystem poisoning, scheduler-induced log and co-location channels, tool/module/MCP poisoning, cross-project data leakage, and coordinated multi-agent exfiltration (Section~\ref{sec:taxonomy}).
  \item A \textbf{defense-gap analysis} showing which existing HPC and agent-security controls detect, partially detect, or miss each attack surface (Section~\ref{sec:defense-gap}).
  \item A \textbf{research agenda}, including the security requirements a trustworthy HPC-agent platform would need to satisfy, and a pointer to an empirical benchmark (HPC-AgentSec) now under development to measure the problem this paper defines (Section~\ref{sec:research-agenda}).
\end{enumerate}

\subsection{Scope and non-goals}
\label{sec:scope}

This paper is a threat characterization and research agenda, not an empirical evaluation. It does not report attack success rates, does not claim a specific model or agent framework is unsafe, and does not propose a complete defense. It also does not attempt to cover model-training-time poisoning, weight extraction, GPU microarchitectural side channels, kernel/hypervisor compromise, or general content-safety jailbreaking --- these are real concerns but orthogonal to the specific claim here, which is about environment-mediated redirection of an agent that already holds valid credentials.
